\clearpage
\section{Contract Selection, Features, and Risk Constraints}\label{app:a}

This appendix separates contract selection, the inputs to each strategy's sizing model, and the admission rules used in the funded contract-level replays. Allocation-state inputs are defined separately in Section~\ref{sec:allocation}.

\subsection{Data and Contract Selection}\label{app:contracts}

The data universe and coverage limitations are described in Section~\ref{sec:data}. Table~\ref{tab:contract-selection} distinguishes the two contract screens.

\begin{table}[H]\centering\small
\caption{Strategy-specific contract selection.}\label{tab:contract-selection}
\setlength{\tabcolsep}{5pt}\renewcommand{\arraystretch}{1.2}
\begin{tabularx}{\textwidth}{@{}>{\hsize=.55\hsize\linewidth=\hsize\raggedright\arraybackslash}X>{\hsize=1.225\hsize\linewidth=\hsize\raggedright\arraybackslash}X>{\hsize=1.225\hsize\linewidth=\hsize\raggedright\arraybackslash}X@{}}\toprule
\textbf{Criterion} & \textbf{Put pair} & \textbf{Straddle}\\\midrule
Expiry & Nearest 30 days; entry requires 7--55 days remaining & ATM-call signal screen requires 21--45 days from next-session entry to expiry\\[4pt]
Strike or delta & Put nearest absolute delta 0.25, within 0.075 & ATM-call strike within 10\% of spot\\[4pt]
Position legs & Selected put on each related underlying & Call and put match on underlying, strike, expiry and adjustment\\
\bottomrule\end{tabularx}\end{table}

Both screens require standard 100-share contracts, positive noncrossed quotes, open interest of at least 100 and quote width no more than 25\% of midpoint. Put-selection ties prefer narrower quotes, higher open interest and then lower contract identifier.

\paragraph{Quote handling and valuation}
Missing entry quotes reject proposals. Corporate events can force exits. Held marks may use last-known IV with current spot, rates and known dividends. European-style valuation and cash settlement approximate expiry; assignment is not fully simulated.

\subsection{Put-Pair Sizing Inputs}\label{app:features}

The tree uses 15 numerical inputs and 12 binary availability flags, giving 27 inputs in total. In Table~\ref{tab:put-sizing-inputs}, starred inputs are required; every other numerical input has an availability flag.

\begingroup\small
\setlength{\tabcolsep}{4pt}\renewcommand{\arraystretch}{1.15}
\begin{longtable}{@{}>{\raggedright\arraybackslash}p{0.32\textwidth}>{\raggedright\arraybackslash}p{\dimexpr0.68\textwidth-8pt\relax}@{}}
\caption{Numerical inputs to put-pair sizing; starred inputs are required.}\label{tab:put-sizing-inputs}\\
\toprule\textbf{Input} & \textbf{Definition}\\\midrule
\endfirsthead
\multicolumn{2}{l}{\textit{Put-pair sizing inputs (continued)}}\\
\toprule\textbf{Input} & \textbf{Definition}\\\midrule
\endhead
\midrule\multicolumn{2}{r}{\textit{Continued on next page}}\\\endfoot
\bottomrule\endlastfoot
\multicolumn{2}{l}{\textit{Signal and volatility}}\\*[3pt]
Mean IV\textsuperscript{*} & Mean implied volatility of the selected puts\\[3pt]
Directional IV gap\textsuperscript{*} & IV gap oriented to the proposed trade direction\\[3pt]
Spread-rank strength\textsuperscript{*} & Strength of the historical IV-spread rank signal\\[3pt]
Directional variance gap & Oriented IV-squared minus trailing-variance gap\\[3pt]
Mean realized volatility & Mean trailing stock realized volatility\\[3pt]
\multicolumn{2}{l}{\textit{Pair relatedness}}\\*[3pt]
Return beta & 126-session stock-return beta\\[3pt]
Return correlation & 126-session stock-return correlation\\[3pt]
\multicolumn{2}{l}{\textit{Contract terms and costs}}\\*[3pt]
Mean time to expiry & Mean days remaining to expiry\\[3pt]
Directional expiry gap & Expiry difference oriented to trade direction\\[3pt]
Commission relative to vega & Training commission expressed in IV points\\[3pt]
\multicolumn{2}{l}{\textit{Market conditions}}\\*[3pt]
VIX & Strictly prior VIX close\\[3pt]
\multicolumn{2}{l}{\textit{Completed directional outcomes}}\\*[3pt]
91-day profitability & Mean completed profit per initial vega over 91 calendar days\\[3pt]
91-day observation count & Number of completed observations in that window\\[3pt]
365-day profitability & Mean completed profit per initial vega over 365 calendar days\\[3pt]
365-day observation count & Number of completed observations in that window\\[3pt]
\end{longtable}\endgroup

An availability flag is one when its numerical field is observed and zero otherwise. Missing optional values are set to zero with flag zero; an observed numerical zero retains flag one. Put trailing variance is 252 times mean squared simple stock returns over 20 sessions, requiring at least ten observations.

\subsection{Straddle Sizing Inputs}\label{app:straddle-features}

The ridge model uses the 14 inputs in Table~\ref{tab:straddle-sizing-inputs}. Benchmark features come from a separate fixed-half straddle portfolio constructed using information available at each date, not the learner's future inventory.

\begingroup\small
\setlength{\tabcolsep}{4pt}\renewcommand{\arraystretch}{1.15}
\begin{longtable}{@{}>{\raggedright\arraybackslash}p{0.34\textwidth}>{\raggedright\arraybackslash}p{\dimexpr0.66\textwidth-8pt\relax}@{}}
\caption{All 14 inputs to the retained straddle ridge model.}\label{tab:straddle-sizing-inputs}\\
\toprule\textbf{Input} & \textbf{Definition}\\\midrule
\endfirsthead
\multicolumn{2}{l}{\textit{Straddle sizing inputs (continued)}}\\
\toprule\textbf{Input} & \textbf{Definition}\\\midrule
\endhead
\midrule\multicolumn{2}{r}{\textit{Continued on next page}}\\\endfoot
\bottomrule\endlastfoot
\multicolumn{2}{l}{\textit{Volatility opportunity and recent changes}}\\[3pt]
Completed RV--IV memory & Historical completed realized-minus-implied volatility memory from Section~\ref{sec:5}\\[3pt]
IV & Current ATM-call implied volatility\\[3pt]
VIX & VIX level available at the proposal date\\[3pt]
Normalized memory & Negative RV--IV memory divided by current IV, with IV floored at 0.01\\[3pt]
Variance ratio & Forecast annualized variance divided by squared IV\\[3pt]
21-session VIX change & Proportional change in VIX over 21 sessions\\[3pt]
5-session IV change & Proportional change in IV over five sessions\\[3pt]
\multicolumn{2}{l}{\textit{Fixed-half benchmark portfolio}}\\[3pt]
21-session Sharpe & Annualized excess Sharpe over 21 sessions\\[3pt]
63-session Sharpe & Annualized excess Sharpe over 63 sessions\\[3pt]
126-session Sharpe & Annualized excess Sharpe over 126 sessions\\[3pt]
63-session volatility & Annualized daily excess-return volatility over 63 sessions\\[3pt]
63-session drawdown & Cash-relative wealth divided by its trailing 63-session maximum, minus one\\[3pt]
Reserve use & Capital reserve divided by current portfolio NAV\\[3pt]
Benchmark availability & Binary flag indicating that the benchmark features are available\\[3pt]
\end{longtable}\endgroup

\subsection{Risk Limits and Admission}\label{app:risk}

The limits in Table~\ref{tab:admission-risk-limits} govern funded contract-level admission, not the outer strategy-index allocation.

\begin{table}[H]\centering\small
\caption{Risk limits used in contract-level admission checks.}\label{tab:admission-risk-limits}
\setlength{\tabcolsep}{5pt}\renewcommand{\arraystretch}{1.25}
\begin{tabularx}{\textwidth}{@{}>{\hsize=.65\hsize\linewidth=\hsize\raggedright\arraybackslash}X>{\hsize=1.35\hsize\linewidth=\hsize\raggedright\arraybackslash}X@{}}\toprule
\textbf{Check} & \textbf{Limit or definition}\\\midrule
Stress loss & At most 5\% of the applicable budget\\[3pt]
Capital reserve & At most 80\% of the applicable budget\\[3pt]
Hedge notional & At most twice the applicable budget\\[3pt]
Gross vega & At most \$20,000 per IV point per \$1 million of reference capital\\[3pt]
Reserve measurement & Gross option premium plus half the netted hedge notional\\[3pt]
Budget scope & Global checks use equity; pair checks use a \$1 million reference scaled by current divided by initial equity\\
\bottomrule\end{tabularx}\end{table}

Stress testing takes each name's worst modeled option-plus-hedge loss across spot shocks of \(-20\%,0,+20\%\) and IV shocks of \(-10,0,+10\) points. These losses are summed without pair-correlation credit. Quote eligibility, inventory and the limits above jointly determine whether a proposed size is admitted.

\FloatBarrier

\subsection{Data and Code Availability}\label{app:reproduction}

Supporting materials, saved results and figures are available at \url{https://github.com/kreitzz/rl-atm-vega-pairs}. Access to the underlying WRDS data requires the relevant subscriptions and vendor permissions; redistribution restrictions apply. Model fitting requires the licensed source cache or suitable archived derived inputs.

Public snapshots support reproduction of the attribution and associated figures from saved derived data, without raw option chains. This verifies archived results rather than independently refitting the models or regenerating all trading experiments. Source locations, verification details and reproduction guidance are provided in the paper's repository README.
\FloatBarrier
